\chapter{Charakterisierungen, Klassen und Partitionen}
\label{b11guide:ch:classes}

\section{Charakterisierungen und Randfälle}

\begin{definition}[Diagonale und Allrelation]
Für eine Menge $A$ seien
\[
 \Delta_A=\{(x,x)\mid x\in A\},\qquad \nabla_A=A\times A.
\]
Sie heißen \emph{Gleichheitsrelation} und \emph{Allrelation} auf $A$.
\end{definition}

\begin{theorem}[Grundmodelle und leere Menge]
\label{b11guide:thm:extremes}
Für jede Menge $A$ gilt
\[
 \Delta_A,\nabla_A\in\EqSet(A),\qquad
 E\in\EqSet(A)\ \Rightarrow\ \Delta_A\subseteq E\subseteq\nabla_A.
\]
Außerdem gilt
\[
\begin{aligned}
 \EqSet(\varnothing)&=\{\varnothing\},&
 \varnothing/\varnothing&=\varnothing,\\
 A/\Delta_A&=\{\{x\}\mid x\in A\},&
 A/\nabla_A&=
 \begin{cases}\{A\},&A\ne\varnothing,\\
                \varnothing,&A=\varnothing.
 \end{cases}
\end{aligned}
\]
Für $E\in\EqSet(A)$ ist $E=\varnothing$ genau dann, wenn
$A=\varnothing$; ebenso ist $A/E=\varnothing$ genau dann,
wenn $A=\varnothing$.
\end{theorem}

\begin{theorem}[Reflexivität und gemeinsamer Ausgangspunkt]
\label{b11guide:thm:euclidean}
Eine Relation $R\subseteq A\times A$ ist genau dann eine
Äquivalenzrelation auf $A$, wenn beide Bedingungen gelten:
\[
\begin{aligned}
 &\forall x\in A\;((x,x)\in R),\\
 &\forall x,y,z\in A\;
   \bigl(((x,y)\in R\land(x,z)\in R)\Rightarrow(y,z)\in R\bigr).
\end{aligned}
\]
Die zweite Bedingung ersetzt Symmetrie und Transitivität durch
eine gemeinsame Bedingung; die Reflexivität bleibt erforderlich.
\end{theorem}

\begin{theorem}[Klassenkriterium für eine beliebige Relation]
\label{b11guide:thm:rowcriterion}
Seien $R\subseteq A\times A$ und
$R_x=\{y\in A\mid(x,y)\in R\}$ für $x\in A$.
Dann sind äquivalent:
\[
\begin{aligned}
 &R\in\EqSet(A),\\
 &\forall x\in A\;(x\in R_x)
 \quad\land\quad
 \forall x,y\in A\;(y\in R_x\Rightarrow R_y=R_x).
\end{aligned}
\]
\end{theorem}

\section{Die grundlegenden Klassensätze}

\begin{theorem}[Klassenzugehörigkeit und Klassengleichheit]
\label{b11guide:thm:class}
Für $E\in\EqSet(A)$ und $x,y\in A$ gilt
$x\in[x]_E\subseteq A$. Ferner sind äquivalent:
\[
\begin{aligned}
 &x\sim_E y,\qquad y\in[x]_E,\qquad x\in[y]_E,\\
 &[x]_E=[y]_E,\qquad [x]_E\cap[y]_E\ne\varnothing,\qquad
 [x]_E\subseteq[y]_E.
\end{aligned}
\]
Zwei Äquivalenzklassen sind somit gleich oder disjunkt.
Jedes Element einer Klasse ist selbst ein Repräsentant dieser Klasse.
\end{theorem}

\begin{theorem}[Gleiche Beziehungen zu allen Elementen]
\label{b11guide:thm:substitution}
Für $E\in\EqSet(A)$ und $x,x',y,y'\in A$ gilt
\[
 x\sim_E x',\ y\sim_E y'
 \quad\Rightarrow\quad
 (x\sim_E y\ \Longleftrightarrow\ x'\sim_E y').
\]
Insbesondere sind für $x,y\in A$ äquivalent:
\[
 x\sim_E y
 \quad\Longleftrightarrow\quad
 \forall z\in A\;(x\sim_E z\ \Longleftrightarrow\ y\sim_E z).
\]
\end{theorem}

\begin{theorem}[Zerlegungssatz]\label{b11guide:thm:partition}
Ist $E\in\EqSet(A)$, so ist $A/E$ eine Partition von $A$.
Jedes $x\in A$ gehört genau einem Block dieser Partition an,
nämlich $[x]_E$. Die Relation läßt sich aus ihren Klassen
zurückgewinnen:
\[
 E=\bigcup_{C\in A/E}(C\times C).
\]
\end{theorem}

\Needspace{19\baselineskip}
\section{Von Partitionen zu Äquivalenzrelationen}

\begin{definition}[Von einer Partition bestimmte Relation]
Für eine Partition $\mathcal C$ von $A$ sei
\[
 E_{\mathcal C}=\bigcup_{C\in\mathcal C}(C\times C).
\]
Für $x,y\in A$ bedeutet $x\sim_{E_{\mathcal C}}y$ also,
daß $x$ und $y$ in demselben Block liegen.
\end{definition}

\begin{theorem}[Partitionssatz und Eindeutigkeit]
\label{b11guide:thm:partitionconverse}
Für jede Partition $\mathcal C$ von $A$ gilt
\[
 E_{\mathcal C}\in\EqSet(A),\qquad A/E_{\mathcal C}=\mathcal C.
\]
Für jede Äquivalenzrelation $E$ auf $A$ gilt umgekehrt
\[
 E_{A/E}=E.
\]
Jede Partition bestimmt damit genau eine Äquivalenzrelation mit
diesen Klassen, und jede Äquivalenzrelation genau ihre Klassenpartition.
Insbesondere gilt für $E,F\in\EqSet(A)$
\[
 E=F\quad\Longleftrightarrow\quad A/E=A/F.
\]
\end{theorem}

